8.- Escribe una gramática que genere un fragmento del español capaz de producir las oraciones ``el gato duerme'' y ``el perro come la croqueta'', pero no ``duerme el gato'' ni ``el croqueta come gato''. Identifica claramente terminales, no terminales, símbolo inicial y producciones.

\vspace{0.5cm}

\vspace{0.2cm}
\noindent Para evitar generar oraciones como ``duerme el gato'' ó ``el croqueta come gato'', voy a generar una serie de reglas:

\begin{enumerate}
    \item Todas las oraciones deben empezar con un artículo
    \item Deben de concordar el género  y el pronombre así se evita ``el croqueta''
    \item Después de un sustantivo, debe de ir un verbo, cumpliendo así la siguiente estructura: \\
    \hspace*{1cm} pronombre + sustantivo + verbo
    \item Pueden haber dos tipos de verbo
    \begin{itemize}
        \item Verbo seguido del complemento
        \item Verbo sin complemento
    \end{itemize}
\end{enumerate}

\vspace{0.5cm}

\noindent $G = \langle V, T, S, P \rangle$

\vspace{0.2cm}
\noindent $T = \langle \text{el}, \text{la}, \text{gato}, \text{perro}, \text{croqueta}, \text{duerme}, \text{come} \rangle$

\vspace{0.2cm}
\noindent $V = \{ S, \text{Sujeto}_m, \text{Sujeto}_f, A_M, A_F, M, F, V_c, V \}$

\vspace{0.4cm}
\noindent $S$ : Símbolo inicial

\vspace{0.2cm}
\noindent $\text{Suj}_m$ : Sujeto masculino

\vspace{0.2cm}
\noindent $\text{Suj}_f$ : Sujeto femenino

\vspace{0.2cm}
\noindent $A_M$ : Artículo masculino

\vspace{0.2cm}
\noindent $A_F$ : Artículo femenino

\vspace{0.2cm}
\noindent $M$ : Sustantivo Masculino

\vspace{0.2cm}
\noindent $F$ : Sustantivo Femenino

\vspace{0.2cm}
\noindent $V_c$ : Verbo con complemento

\vspace{0.2cm}
\noindent $V$ : Verbo sin complemento

\vspace{0.6cm}
\underline{\noindent Reglas de producción}

\vspace{0.3cm}
\noindent $S \rightarrow \text{Suj}_m V \mid \text{Suj}_m V_c \, \text{Suj}_f$

\vspace{0.25cm}
\noindent $\text{Suj}_m \rightarrow A_M M$

\vspace{0.25cm}
\noindent $A_M \rightarrow \text{el}$

\vspace{0.25cm}
\noindent $A_F \rightarrow \text{la}$

\vspace{0.25cm}
\noindent $M \rightarrow \text{gato} \mid \text{perro}$

\vspace{0.25cm}
\noindent $F \rightarrow \text{croqueta}$

\vspace{0.25cm}
\noindent $V \rightarrow \text{duerme}$

\vspace{0.25cm}
\noindent $V_c \rightarrow \text{come}$

\vspace{0.3cm}
\noindent ``el gato duerme''

\vspace{0.25cm}
\noindent $S \rightarrow \text{Suj}_m V \rightarrow A_M M V \rightarrow \text{el} \, M V$ \\
\hspace*{0.5cm} $\rightarrow \text{el gato} \, V \rightarrow \text{el gato duerme}$

\vspace{0.6cm}
\noindent ``el perro come la croqueta''

\vspace{0.25cm}
\noindent $S \rightarrow \text{Suj}_m V_c \, \text{Suj}_f \rightarrow A_M M V_c \, \text{Suj}_f$ \\
\hspace*{0.5cm} $\rightarrow \text{el perro come} \, \text{Suj}_f \rightarrow \text{el perro come} \, A_F F$ \\
\hspace*{0.5cm} $\rightarrow \text{el perro come la} \, F \rightarrow \text{el perro come la croqueta}$